\documentclass[10pt,a4paper]{amsart}
\usepackage[margin=19mm]{geometry}
\usepackage{amsmath,amssymb,mathtools}
\usepackage[scr=rsfs,cal=euler]{mathalfa}
\usepackage{xfrac}
\usepackage[unicode,hidelinks,bookmarks=false]{hyperref}
\newcommand{\ie}{i.e.\ }
\newcommand{\cf}{cf.\ }
\newcommand{\resp}{respectively\ }
\newcommand{\Subsection}{\subsection}
\providecommand{\nobreakdash}{\nobreak\hbox{-}}
\setlength{\emergencystretch}{2em}
\multlinegap0pt
\usepackage{bm}
\usepackage{bbm}
\usepackage{dsfont}
\usepackage{xspace}
\usepackage{cancel}
\usepackage{mathtools}
\usepackage{amsthm}
\makeatletter
\@ifundefined{theorem}{\newtheorem{theorem}{Theorem}}{}
\@ifundefined{assumption}{\newtheorem{assumption}[theorem]{Assumption}}{}
\makeatother
\title[Assortment Optimization under the Multinomial Logit Model wi]{Assortment Optimization under the Multinomial Logit Model with Covering Constraints}
\author{Omar El Housni, Qing Feng, Huseyin Topaloglu}
\date{Source: arXiv:2411.10310v2 (CC BY 4.0)}
\begin{document}
\maketitle

\noindent\emph{Source and licence.} Assortment Optimization under the Multinomial Logit Model with Covering Constraints. Version arXiv:2411.10310v2 (CC BY 4.0), distributed under the Creative Commons Attribution 4.0 International licence, as recorded in the arXiv licence metadata for this exact version and shown on the abstract page https://arxiv.org/abs/2411.10310v2. Full rights notice: licenses/arxiv-2411.10310-CC-BY-4.0.txt.\par\smallskip

\noindent\textit{Editorial scope.} Counted unit: the Theorem whose source label is \texttt{None} in the article (source0.tex lines 2010--2012, source label \texttt{None}), delivered together with the article's own complete original proof of it (source0.tex lines 2014--2109, one \texttt{proof} environment of 96 lines). No mathematics is written, repaired or re-derived: statement and proof are the article's own text line for line. Required context retained uncounted: prob:deterministic\_single\_segment (lines 362--369); prob:framing (lines 1976--1982); assump:monotone\_browsing\_distribution (lines 1988--1990); assump:enough\_position (lines 1992--1994); eq:best\_candidate\_solution (lines 2006--2008). Every definition, display and label of those lines is retained unchanged. Omitted from this package and not redistributed: 1--361, 370--1975, 1983--1987, 1991--1991, 1995--2005, 2009--2009, 2013--2013, 2110--2118. Nothing inside a retained excerpt is omitted or altered: every retained body is delivered exactly as the article's lines are written, so no omission receipt of any kind accompanies this package. Citations: the retained excerpts carry no citation command at all, so no bibliography entry of the article is transcribed and no reference is redistributed. Reference adaptation: none was needed and none was made; every reference inside the delivered unit resolves inside it, so no unresolved reference or citation is produced. Printed slips kept literal: no printed word, symbol, hypothesis or number of the counted statement or of its proof was altered. Layout and class: the article's own class is \texttt{informs3} with packages including endnotes, algorithm, algorithmicx, algpseudocode, amsfonts, subfigure, optidef, graphicx, hyperref, comment, appendix, cleveref, tikz, multirow; the wrapper is the lane's pinned \texttt{amsart} editorial wrapper at 10pt on A4, which restates the theorem-like environments the retained text needs and adds the article's own macro definitions that the retained text needs (none), with extra wrapper packages (bm, bbm, dsfont, xspace, cancel, mathtools). The delivered numbering is the unit's own. Assets: no retained span contains a figure environment, a graphics-inclusion command, a PDF-page inclusion, a table or a TikZ picture; no image file is copied into this package, the package declares no asset dependency and no image file is redistributed with it. Licence: version arXiv:2411.10310v2 of this article is distributed under the Creative Commons Attribution 4.0 International licence, as recorded in the arXiv licence metadata for this exact version and shown on the abstract page https://arxiv.org/abs/2411.10310v2. A full rights notice is delivered at licenses/arxiv-2411.10310-CC-BY-4.0.txt. Characters: every retained excerpt is pure ASCII, so the delivered engine, XeLaTeX with its default Latin Modern text font, reports no missing character.
\begin{equation}
\label{prob:deterministic_single_segment}
\tag{DAOC}
\begin{aligned}
&\max_{S\subseteq\mathcal{N}} && R(S)\\
&\textup{s.t.} && |S\cap C_k|\geq \ell_k,\ \forall k\in\{1,2,\dots,K\}.\\
\end{aligned}
\end{equation}

\begin{equation}
\label{prob:framing}
\begin{aligned}
&\max_{X:\mathcal{G}\to\mathcal{N}\cup\{0\}} && \mathcal{R}(X)\\
&\text{s.t.} && |X(1:G)\cap C_k|\geq \ell_k,\ \forall k\in\{1,2,\dots,K\}.
\end{aligned}
\end{equation}

\begin{assumption}\label{assump:monotone_browsing_distribution}
The browsing probabilities satisfy $\beta_g\leq \beta_{g'}$ for all $g\geq g'$.
\end{assumption}

\begin{assumption}\label{assump:enough_position}
The total number of positions $G$ satisfies $G\geq 3n$.
\end{assumption}

\begin{equation}\label{eq:best_candidate_solution}
Y=\underset{X\in\{X_i\}_{i=1}^n}{\mathrm{argmax}}\,\mathcal{R}(X).
\end{equation}

\begin{theorem}
Under Assumptions~\ref{assump:monotone_browsing_distribution}~and~\ref{assump:enough_position}, $Y$ defined in \eqref{eq:best_candidate_solution} provides a $1/2\log_2(8n)$-approximation to Problem~\eqref{prob:framing}.
\end{theorem}

\begin{proof}

Let $X^*:\mathcal{G}\to\mathcal{N}\cup\{0\}$ be an optimal solution and $\text{OPT}$ be the optimal value of Problem~\eqref{prob:framing}. We define $\hat{X}$ as
\begin{equation*}
\hat{X}(g)=\begin{cases}
X^*(g),\ &\text{if}\ X^*(i)\notin X^*(1:g-1);\\
0,\ &\text{otherwise}.
\end{cases}
\end{equation*}
In other words, in the definition of $\hat{X}$, we remove duplicate products by only keeping each products in the position where it first appears. By definition of $\hat{X}$ we have $i\in \hat{X}(1:g)$ implies $i\in X^*(1:g)$. Furthermore, if $i\in X^*(1:g)$, then the first position where product $i$ appears is among positions $1$ through $g$, thus $i\in \hat{X}(1:g)$. Therefore we conclude that $X^*(1:g)=\hat{X}(1:g)$ for all $g\in\{1,2,\dots,G\}$, thus $\text{OPT}=\mathcal{R}(X^*)=\mathcal{R}(\hat{X})$. We further have
\begin{equation}\label{eq:framing_1}
\begin{split}
\text{OPT}&=\sum_{g=1}^G\beta_gR(\hat{X}(1:g))=\sum_{g=1}^G\beta_g\sum_{i\in \hat{X}(1:g)}r_i\phi(i,\hat{X}(1:g))\\
&=\sum_{g=1}^G\beta_g\sum_{g'=1}^g r_{\hat{X}(g')}\phi(\hat{X}(g'),\hat{X}(1:g))\\
&=\sum_{g'=1}^G\sum_{g=g'}^G \beta_g r_{\hat{X}(g')}\phi(\hat{X}(g'),\hat{X}(1:g)).
\end{split}
\end{equation}
Here for notational brevity we define $r_0=0$, and the third equality holds because by definition of $\hat{X}$, every product appears at most once in the product framing $\hat{X}$. We set $\ell_{\max}=\lceil\log_2(n)\rceil$, and define
\begin{equation*}
g_0=0,\ g_\ell=2^{\ell-1},\ \forall \ell\in\{1,2,\dots,\ell_{\max}\},\ g_{\ell_{\max}+1}=n,\ g_{\ell_{\max}+2}=G.
\end{equation*}

For any $\ell\in\{0,1,\dots,\ell_{\max}+1\}$, we define $\hat{X}_\ell$ as the product framing where positions $g_\ell+1$ through $g_{\ell+1}$ are kept same as $\hat{X}$ while all other positions are left empty, i.e.,
\begin{equation*}
\hat{X}_\ell(g)=\begin{cases}
\hat{X}(g),\ &\text{if}\ g_{\ell}+1\leq g\leq g_{\ell+1},\\
0,\ &\text{otherwise}.
\end{cases}
\end{equation*}
Then by~\eqref{eq:framing_1}~we get
\begin{align*}
\text{OPT}=&\sum_{g'=1}^G\sum_{g=g'}^G \beta_g r_{\hat{X}(g')}\phi(\hat{X}(g'),\hat{X}(1:g))=\sum_{\ell=0}^{\ell_{\max}+1}\sum_{g'=g_{\ell}+1}^{g_{\ell+1}}\sum_{g=g'}^G \beta_g r_{\hat{X}(g')}\phi(\hat{X}(g'),\hat{X}(1:g))\\
\leq&\sum_{\ell=0}^{\ell_{\max}+1}\sum_{g'=g_\ell+1}^{g_{\ell+1}}\sum_{g=g'}^G\beta_g r_{\hat{X}(g')}\phi(\hat{X}(g'),\hat{X}(g_\ell+1:\min\{g,g_{\ell+1}\}))=\sum_{\ell=0}^{\ell_{\max}+1}\mathcal{R}(\hat{X}_\ell).
\end{align*}
Here the inequality is because under the MNL model, $\phi(i,S)\geq \phi(i,S')$ if $i\in S$ and $S\subseteq S'$. We define $\ell^*=\mathrm{argmax}_{\ell}\,\mathcal{R}(\hat{X}_\ell)$, then we have
\begin{equation*}
\mathcal{R}(\hat{X}_{\ell^*})\geq \dfrac{1}{\ell_{\max}+2}\sum_{\ell=0}^{\ell_{\max}+1}\mathcal{R}(\hat{X}_\ell)\geq \dfrac{\text{OPT}}{\ell_{\max}+2}.
\end{equation*}

Then we prove that there exists $i\in\{1,2,\dots,n\}$ such that $\mathcal{R}(X_i)\geq \mathcal{R}(\hat{X}_{\ell^*})/2$. If $\ell^*\neq 0$ and $\ell^*\neq \ell_{\max}+1$, we set $i=g_{\ell^*+1}-g_{\ell^*}$, then
\begin{equation*}
\mathcal{R}(\hat{X}_{\ell^*})=\sum_{g=g_{\ell^*}+1}^{G}\beta_g R(\hat{X}(g_{\ell^*}+1:\min\{g,g_{\ell^*+1}\}))\leq \sum_{g=g_{\ell^*}+1}^{G}\beta_g R(S_i^*)\leq \sum_{g=i}^G\beta_gR(S_i^*).
\end{equation*}
Here the first inequality is due to $|\hat{X}(g_{\ell^*}+1:\min\{g,g_{\ell^*+1}\})|\leq g_{\ell^*+1}-g_{\ell^*}=i$, thus by definition of $S_i^*$ we have $R(S_i^*)\geq R(\hat{X}(g_{\ell^*}+1:g))$ for all $g_{\ell^*}+1\leq g\leq G$. The second inequality is due to the assumption that $\ell^*\neq \ell_{\max}+1$ and $\ell^*\neq 0$, thus $i=g_{\ell^*+1}-g_{\ell^*}\leq g_{\ell^*}$. Furthermore, since the approximate solution $\hat{S}$ of \eqref{prob:deterministic_single_segment} satisfies $|\hat{S}|\leq n$, under product framing $X_i$, if a customer browses up to $g$ for some $i\leq g\leq G-n$, the set of products offered to the customer is $S_i^*$. Therefore
\begin{equation*}
\mathcal{R}(X_i)\geq \sum_{g=i}^{G-n}\beta_g R(S_i^*).
\end{equation*}
We further have
\begin{equation*}
\sum_{g=G-n+1}^G \beta_g\leq \sum_{g=G-2n+1}^{G-n}\beta_g\leq \sum_{g=i}^{G-n}\beta_g.
\end{equation*}
Here the first inequality is because by Assumption~\ref{assump:monotone_browsing_distribution}~we get $\beta_g\leq \beta_{g-n}$ for all $G-n+1\leq g\leq G$. The second inequality is because we have assumed $\ell^*\neq \ell_{\max}+1$, thus $i=g_{\ell^*+1}-g_{\ell^*}\leq g_{\ell^*+1}\leq n$, and by Assumption~\ref{assump:enough_position}~we have $G-2n+1\geq n\geq i$. Therefore we have $\sum_{g=i}^{G-n}\beta_g\leq \sum_{g=i}^G\beta_g/2$. Then we get
\begin{equation*}
\mathcal{R}(X_i)\geq \sum_{g=i}^{G-n}\beta_g R(S_i^*)\geq \dfrac{1}{2}\sum_{g=i}^G \beta_gR(S_i^*)\geq \dfrac{1}{2}\mathcal{R}(\hat{X}_{\ell^*})\geq \dfrac{\text{OPT}}{2(\ell_{\max}+2)}.
\end{equation*}

If $\ell^*=0$, since by definition $g_1=1$, we have
\begin{equation*}
\mathcal{R}(\hat{X}_0)=\sum_{g=1}^G \beta_g R(\{\hat{X}_0(1)\})=R(\{\hat{X}(1)\})\cdot\sum_{g=1}^G\beta_g=R(\{\hat{X}(1)\})\leq R(S_1^*).
\end{equation*}
On the other hand, we have
\begin{equation*}
\mathcal{R}(X_1)\geq \sum_{g=1}^{G-n}\beta_gR(S_1^*).
\end{equation*}
By Assumption~\ref{assump:monotone_browsing_distribution}~we also have
\begin{equation*}
\sum_{g=G-n+1}^G\beta_g\leq \sum_{g=G-2n+1}^{G-n}\beta_g\leq \sum_{g=1}^{G-n}\beta_g.
\end{equation*}
Therefore we have $\sum_{g=1}^{G-n}\beta_g\geq \sum_{g=1}^G\beta_g/2$, therefore
\begin{equation*}
\mathcal{R}(X_1)\geq \sum_{g=1}^{G-n}\beta_g R(S_1^*)\geq \dfrac{1}{2}\sum_{g=1}^G \beta_gR(S_1^*)\geq \mathcal{R}(\hat{X}_0)\geq \dfrac{\text{OPT}}{2(\ell_{\max}+2)}.
\end{equation*}

If $\ell^*=\ell_{\max}+1$, then we have
\begin{equation*}
\mathcal{R}(\hat{X}_{\ell^*})=\sum_{g=n+1}^G \beta_gR(\hat{X}(n:g))\leq \sum_{g=n+1}^G\beta_gR(S_n^*).
\end{equation*}
Here the inequality is because $S_n^*$ is the unconstrained optimal assortment, thus $R(S_n^*)\geq R(S)$ for all $S\subseteq\mathcal{N}$. On the other hand, we have
\begin{equation*}
\mathcal{R}(X_n)\geq \sum_{g=n}^{G-n}\beta_g R(S_n^*).
\end{equation*}
We also have
\begin{equation*}
\sum_{g=G-n+1}^{G}\beta_g\leq \sum_{g=G-2n+1}^{G-n}\beta_g\leq \sum_{g=n}^{G-n}\beta_g.
\end{equation*}
Here the first inequality is because by Assumption~\ref{assump:monotone_browsing_distribution}~we get $\beta_g\leq \beta_{g-n}$ for all $G-n+1\leq g\leq G$. The second inequality is because by Assumption~\ref{assump:enough_position} we have $G-2n+1\geq n$. Therefore we have $\sum_{g=n}^{G-n}\beta_g\geq \sum_{g=n}^G \beta_g/2$, thus
\begin{equation*}
\mathcal{R}(X_n)\geq \sum_{g=n}^{G-n}\beta_g R(S_n^*)\geq \dfrac{1}{2}\sum_{g=n+1}^G \beta_gR(S_n^*)\geq \dfrac{1}{2}\mathcal{R}(\hat{X}_{\ell^*})\geq \dfrac{\text{OPT}}{2(\ell_{\max}+2)}.
\end{equation*}
Therefore in all three cases, there exists $i\in\{1,2,\dots,n\}$ such that $\mathcal{R}(X_i)\geq \text{OPT}/2(\ell_{\max}+2)$. Since we define $Y=\mathrm{argmax}_{X\in\{X_i\}_{i=1}^n}\mathcal{R}(X)$, we have
\begin{equation*}
\mathcal{R}(Y)\geq \dfrac{\text{OPT}}{2(\ell_{\max}+2)}\geq \dfrac{\text{OPT}}{2\log_2(8n)}.
\end{equation*}

Finally we show that for any $i\in\{1,2,\dots,n\}$, $X_i$ satisfies the covering constraints. Since $\hat{S}$ is always offered at the last $|\hat{S}|$ positions of $X_i$, we get $X_i(1:G)\supseteq \hat{S}$, thus $|X_i(1:G)\cap C_k|\geq |\hat{S}\cap C_k|\geq \ell_k$ for all $k\in\{1,2,\dots,K\}$. Therefore $Y$ satisfies the covering constraints, thus is a $1/2\log(8n)$-approximation to Problem~\eqref{prob:framing}.
\end{proof}

\end{document}
